\section{Results}
\label{sec:results}

\subsection{Where land value sits}
\label{sec:whereland}

The distributional consequences of an LVT are determined by the distribution of its base, which we describe first. Table~\ref{tab:land_income} reports average land value, property wealth and land shares by decile of equivalised household net income. Two features stand out. First, average land value rises with income, but weakly and unevenly: it is \pounds 95{,}302 in the poorest decile and \pounds 442{,}935 in the richest, a ratio of 4.6 to 1 against an income ratio of 6.5 to 1, and the profile is far from monotonic --- decile 2 (\pounds 217{,}644) holds nearly twice as much land on average as decile 3 (\pounds 115{,}702). Some of this reflects genuine heterogeneity, principally asset-rich retired households with low current income; some reflects imputation noise in a base constructed from survey property wealth. Either way, the income gradient of the tax base is shallow, which is what drives the results in Section~\ref{sec:poverty}. Second, the top two income deciles together hold 35.5 per cent of land, whereas, as Section~\ref{sec:wealth_impact} shows, the top two \emph{wealth} deciles hold 66.6 per cent.

Table~\ref{tab:land_income} separates the two components of the base, which are not comparable quantities: household land is a share of the household's own property wealth, while allocated corporate land is an attributed claim on corporate balance sheets that households holding no property may still bear. Reporting them together against property wealth --- as a single ``land value'' column would --- can make land appear to exceed the property wealth it is derived from.

\begin{table}[htbp]
\centering
\caption{Land value and property wealth by decile of equivalised household net income, 2026--27}
\label{tab:land_income}
\begin{tabular}{lrrrrrr}
\toprule
Income & Avg.\ income & Avg.\ household & Avg.\ corporate & Avg.\ total & Avg.\ property & Share of \\
decile & (\pounds/yr) & land (\pounds) & land (\pounds) & land (\pounds) & wealth (\pounds) & land (\%) \\
\midrule
1  & 21{,}701  & 75{,}848  & 19{,}454  & 95{,}302  & 106{,}712 & 4.7 \\
2  & 31{,}068  & 171{,}344 & 46{,}300  & 217{,}644 & 198{,}309 & 8.8 \\
3  & 37{,}182  & 89{,}340  & 26{,}362  & 115{,}702 & 144{,}651 & 5.0 \\
4  & 38{,}400  & 152{,}099 & 50{,}164  & 202{,}263 & 204{,}378 & 9.3 \\
5  & 43{,}127  & 144{,}852 & 52{,}305  & 197{,}157 & 213{,}009 & 9.0 \\
6  & 53{,}985  & 160{,}086 & 50{,}458  & 210{,}544 & 231{,}152 & 8.5 \\
7  & 55{,}357  & 138{,}281 & 48{,}129  & 186{,}410 & 221{,}029 & 8.1 \\
8  & 69{,}125  & 191{,}548 & 88{,}673  & 280{,}220 & 308{,}960 & 10.9 \\
9  & 77{,}096  & 294{,}523 & 127{,}118 & 421{,}641 & 451{,}521 & 17.5 \\
10 & 141{,}327 & 294{,}942 & 147{,}992 & 442{,}935 & 429{,}497 & 18.0 \\
\bottomrule
\end{tabular}

\medskip
{\footnotesize\emph{Source:} Authors' calculations using PolicyEngine UK.}
\end{table}

Regional variation is at least as pronounced as variation by income. Average household land value is \pounds 369{,}802 in London and \pounds 286{,}311 in the South East, compared with \pounds 153{,}631 in the North East and \pounds 159{,}535 in Scotland (Figure~\ref{fig:land_region}) --- a ratio of 2.4 to 1 that reflects both higher property prices and higher land shares of those prices in the South. Across household types, pensioner couples hold the most land on average (\pounds 431{,}377), followed by working-age couples without children (\pounds 303{,}440); lone parents, who are predominantly renters, hold very little (\pounds 7{,}797) (Figure~\ref{fig:land_famtype}). Ranked by total wealth rather than income, concentration is far greater: the bottom three wealth deciles hold essentially no land, while the top wealth decile averages \pounds 1{,}088{,}892 --- 47.4 per cent of all land (Table~\ref{tab:wealth_impact} below). Overall, 17.2 per cent of households hold no land.

\begin{figure}[htbp]
\centering
\safefig{../results/figures/fig2_land_by_income_decile.png}
\caption{Average land value by income decile.}
\label{fig:land_income}
\end{figure}

\begin{figure}[htbp]
\centering
\safefig{../results/figures/fig3_land_by_region.png}
\caption{Average household land value by region.}
\label{fig:land_region}
\end{figure}

\begin{figure}[htbp]
\centering
\safefig{../results/figures/fig4_land_by_family_type.png}
\caption{Average household land value by family type.}
\label{fig:land_famtype}
\end{figure}

\subsection{Revenue and the rate schedule}
\label{sec:revenue}

Because the base is fixed in the static framework, LVT revenue is linear in the rate: each 0.1 percentage point raises approximately \pounds 7.5 billion. Table~\ref{tab:revenue} reports revenue across the rate grid. A 0.5 per cent rate raises \pounds 37.3 billion --- \pounds 20.3 billion less than council tax --- while a 1 per cent rate raises \pounds 74.6 billion and a 5 per cent rate raises \pounds 373.2 billion. The budget-neutral rate is 0.77 per cent (equation~\ref{eq:rate}). Decomposing the base at a 1 per cent rate, household land contributes \pounds 54.0 billion and corporate land \pounds 20.6 billion of the \pounds 74.6 billion total. The corporate component --- 28 per cent of revenue --- falls on households in proportion to their shareholdings, and its inclusion permits a lower rate for any given revenue target than a purely residential base would require.

\begin{table}[htbp]
\centering
\caption{LVT revenue by rate, 2026--27}
\label{tab:revenue}
\begin{tabular}{lrrrr}
\toprule
Rate & LVT revenue & Council tax & Net revenue & Avg.\ LVT per \\
     & (\pounds bn) & (\pounds bn) & (\pounds bn) & household (\pounds) \\
\midrule
0.5\% & 37.3  & 57.6 & $-$20.3 & 1{,}163 \\
1.0\% & 74.6  & 57.6 & 17.0    & 2{,}327 \\
1.5\% & 111.9 & 57.6 & 54.3    & 3{,}490 \\
2.0\% & 149.3 & 57.6 & 91.6    & 4{,}653 \\
3.0\% & 223.9 & 57.6 & 166.3   & 6{,}980 \\
5.0\% & 373.2 & 57.6 & 315.5   & 11{,}634 \\
\bottomrule
\end{tabular}

\medskip
{\footnotesize\emph{Note:} Net revenue is LVT revenue less replaced council tax revenue. \emph{Source:} Authors' calculations using PolicyEngine UK.}
\end{table}

\begin{figure}[htbp]
\centering
\safefig{../results/figures/fig7_revenue_by_rate.png}
\caption{LVT revenue by rate against the council tax replacement target.}
\label{fig:revenue}
\end{figure}

\subsection{Distributional impact by income decile}
\label{sec:income_impact}

Table~\ref{tab:income_impact} and Figure~\ref{fig:net_income} report the average change in household net income at the budget-neutral 0.77 per cent rate, by income decile. Because the swap generates no benefit interactions --- a result we verify in Section~\ref{sec:robustness} --- the change is exactly the abolished council tax bill less the land tax bill, $\Delta y_i = \mathrm{CT}_i - r^{*}L_i$.

The pattern across income deciles is best described as flat with a step down at the top. The poorest decile gains most in both absolute and relative terms, \pounds 647 per year or 3.0 per cent of net income, because its average council tax bill of \pounds 1{,}383 --- inflated relative to its means by the compressed banding structure --- is nearly twice the \pounds 736 land tax charged on its modest holdings. Deciles 3 and 7 gain \pounds 578 and \pounds 455; deciles 5, 6 and 8 are close to unchanged; and deciles 2 and 4 record small average \emph{losses} of \pounds 163 and \pounds 60, reflecting the uneven land profile documented in Table~\ref{tab:land_income} rather than anything systematic about those income levels. Only in deciles 9 and 10 is the loss substantial and clearly attributable to income: \pounds 1{,}016 and \pounds 963 per year, as land holdings of \pounds 420{,}000--\pounds 443{,}000 attract liabilities of \pounds 3{,}250--\pounds 3{,}420 against council tax savings of \pounds 2{,}240--\pounds 2{,}460. In percentage terms the loss is larger for decile 9 ($-1.3$ per cent) than for decile 10 ($-0.7$ per cent), because top-decile incomes are considerably higher while average land holdings are only slightly larger.

We stress that the middle of this profile is not precisely estimated. The differences between adjacent deciles in the range 2--8 are of the order of a few hundred pounds a year on average net incomes of \pounds 31{,}000--\pounds 69{,}000, and they rest on imputed land values whose decile means are visibly noisy. The result that should be read from Table~\ref{tab:income_impact} is that the reform is roughly distributionally neutral across the bottom eight deciles taken together and clearly costly to the top two --- not that any particular middle decile gains or loses.

\begin{table}[htbp]
\centering
\caption{Impact of the 0.77\% budget-neutral LVT by income decile}
\label{tab:income_impact}
\begin{tabular}{lrrrrrr}
\toprule
Income & Avg.\ LVT & Avg.\ CT & Avg.\ net & \% of net & Winners & Losers \\
decile & (\pounds) & saved (\pounds) & change (\pounds) & income & (\%) & (\%) \\
\midrule
1  & 736     & 1{,}383 & $+$647    & $+$3.0 & 79.0 & 18.3 \\
2  & 1{,}681 & 1{,}518 & $-$163    & $-$0.5 & 72.0 & 26.9 \\
3  & 894     & 1{,}472 & $+$578    & $+$1.6 & 76.5 & 22.4 \\
4  & 1{,}562 & 1{,}502 & $-$60     & $-$0.2 & 72.6 & 26.6 \\
5  & 1{,}523 & 1{,}668 & $+$145    & $+$0.3 & 69.1 & 29.7 \\
6  & 1{,}626 & 1{,}805 & $+$179    & $+$0.3 & 70.9 & 28.6 \\
7  & 1{,}440 & 1{,}895 & $+$455    & $+$0.8 & 76.0 & 23.7 \\
8  & 2{,}164 & 2{,}199 & $+$35     & $+$0.1 & 67.9 & 30.9 \\
9  & 3{,}256 & 2{,}240 & $-$1{,}016 & $-$1.3 & 48.1 & 51.4 \\
10 & 3{,}421 & 2{,}458 & $-$963    & $-$0.7 & 48.6 & 50.8 \\
\bottomrule
\end{tabular}

\medskip
{\footnotesize\emph{Note:} ``CT saved'' denotes the abolished council tax liability net of council tax reduction. Winners and losers do not sum to 100 per cent because a small share of households are unaffected. \emph{Source:} Authors' calculations using PolicyEngine UK.}
\end{table}

\begin{figure}[htbp]
\centering
\safefig{../results/figures/fig1_net_change_by_income_decile.png}
\caption{Average net income change by income decile at the 0.77\% budget-neutral rate.}
\label{fig:net_income}
\end{figure}

\subsection{Distributional impact by wealth decile}
\label{sec:wealth_impact}

Table~\ref{tab:wealth_impact} and Figure~\ref{fig:net_wealth} report the corresponding results when households are ranked by total wealth rather than income.\footnote{We construct equal-weight wealth deciles directly. PolicyEngine UK's own \texttt{household\_wealth\_decile} variable is degenerate at the bottom of the distribution --- the mass of households with zero or negative net wealth is tied, leaving decile one empty and placing 22.6 per cent of households in decile two --- which would make the decile averages in this table incomparable.} The gradient here is monotonic and steep, in sharp contrast to the income ranking. The bottom two wealth deciles own essentially no land and retain their council tax savings in full, gaining \pounds 898 and \pounds 902 per year; deciles 3--7 gain between \pounds 523 and \pounds 1{,}480; decile 8 loses \pounds 268; decile 9 loses \pounds 1{,}069; and the top wealth decile --- averaging \pounds 1.09 million of land --- loses \pounds 5{,}752 per year, or 8.8 per cent of its net income, with an average LVT liability of \pounds 8{,}410 against council tax savings of \pounds 2{,}658. In the top wealth decile 98.4 per cent of households lose; in each of the bottom seven, more than 82 per cent gain. The incidence of the reform tracks land wealth almost mechanically, which is exactly what the instrument is designed to do.\footnote{Revenue neutrality holds in aggregate --- LVT revenue equals abolished council tax revenue at \pounds 57.6 billion --- so the population-weighted average net-income change is zero. It is identical whether households are grouped by income or by wealth decile, since the two tables show the same households re-sorted; the per-decile averages need not sum to zero because decile group sizes are equal in population but not in the underlying quantities.}

\begin{table}[htbp]
\centering
\caption{Impact of the 0.77\% budget-neutral LVT by wealth decile}
\label{tab:wealth_impact}
\begin{tabular}{lrrrrrr}
\toprule
Wealth & Avg.\ land & Avg.\ LVT & Avg.\ CT & Avg.\ net & Winners & Losers \\
decile & value (\pounds) & (\pounds) & saved (\pounds) & change (\pounds) & (\%) & (\%) \\
\midrule
1  & 1             & 0       & 898     & $+$898     & 82.1 & 13.7 \\
2  & 29            & 0       & 902     & $+$902     & 83.2 & 11.4 \\
3  & 1{,}265       & 10      & 1{,}177 & $+$1{,}168 & 91.2 & 8.1 \\
4  & 25{,}263      & 195     & 1{,}675 & $+$1{,}480 & 94.4 & 5.6 \\
5  & 76{,}699      & 592     & 1{,}833 & $+$1{,}241 & 91.8 & 8.2 \\
6  & 138{,}998     & 1{,}074 & 2{,}027 & $+$953     & 89.5 & 10.5 \\
7  & 212{,}347     & 1{,}640 & 2{,}163 & $+$523     & 82.9 & 17.1 \\
8  & 321{,}469     & 2{,}483 & 2{,}215 & $-$268     & 48.0 & 52.0 \\
9  & 451{,}386     & 3{,}486 & 2{,}417 & $-$1{,}069 & 20.0 & 80.0 \\
10 & 1{,}088{,}892 & 8{,}410 & 2{,}658 & $-$5{,}752 & 1.6  & 98.4 \\
\bottomrule
\end{tabular}

\medskip
{\footnotesize\emph{Note:} Deciles of total household wealth. \emph{Source:} Authors' calculations using PolicyEngine UK.}
\end{table}

\begin{figure}[htbp]
\centering
\safefig{../results/figures/fig5_net_change_by_wealth_decile.png}
\caption{Average net income change by wealth decile at the 0.77\% budget-neutral rate.}
\label{fig:net_wealth}
\end{figure}

\subsection{Winners and losers within deciles}
\label{sec:winners}

Decile averages conceal substantial within-decile variation, because land holdings vary widely within any given income band: an outright owner-occupier in London and a private renter in the North East may occupy the same income decile. At the budget-neutral rate, 68.2 per cent of all households gain. The share of gaining households declines only gradually with income through deciles 1--8 (from 79.0 to 67.9 per cent) before falling to just under one half in deciles 9--10 (Figure~\ref{fig:winners}). Even in the top income decile, 48.6 per cent of households gain --- predominantly high earners who rent or own property of modest value --- while 18.3 per cent of bottom-decile households lose, typically low-income owners of valuable land. Sorted by wealth, the gradient is considerably steeper: the share of gaining households exceeds 82 per cent in each of the bottom seven wealth deciles and falls to 1.6 per cent in the top decile. The 17.2 per cent of households holding no land gain \pounds 832 per year on average, and 81.5 per cent of them gain; the remainder are largely households whose council tax was already fully offset by council tax reduction.

That the winner share is high while the average income-decile effects are small is not a contradiction. The reform moves a fixed \pounds 57.6 billion from a base that is broad and shallow (council tax, capped at the top) to one that is narrow and deep (land, highly concentrated). Most households pay a little less; a small minority pay a great deal more. This is visible in the wealth ranking and almost invisible in the income ranking.

\begin{figure}[htbp]
\centering
\safefig{../results/figures/fig6_winners_losers_by_decile.png}
\caption{Share of winners and losers by income decile at the 0.77\% budget-neutral rate.}
\label{fig:winners}
\end{figure}

\subsection{Poverty and inequality}
\label{sec:poverty}

Table~\ref{tab:poverty} reports poverty and inequality across the rate grid, measured against baseline-fixed absolute poverty thresholds. At the budget-neutral 0.77 per cent rate the two poverty measures move in opposite directions. Absolute poverty after housing costs falls from 14.75 to 14.04 per cent, a reduction of 0.70 percentage points, while poverty before housing costs is essentially unchanged, edging up from 11.53 to 11.58 per cent ($+0.04$ points). The Gini coefficient of equivalised household net income is likewise unchanged, rising by 0.0003 from a baseline of 0.3641.

The BHC and AHC divergence is informative rather than a puzzle. The AHC measure nets housing costs out of income, and the households lifted above the AHC line are disproportionately low-income renters, who gain their entire council tax bill and owe almost no land tax. The BHC measure does not, so it gives more weight to low-income owner-occupiers, who lose. The near-zero BHC effect is the sum of these two flows, and it is small enough that we would not claim a poverty reduction on the BHC measure at all: the honest statement is that the revenue-neutral swap reduces after-housing-cost poverty modestly and leaves before-housing-cost poverty unchanged.

The grid shows how sensitive this is to the rate, and therefore to the replacement target. At 0.5 per cent --- which is not revenue-neutral, reducing net revenue by \pounds 20.3 billion --- poverty falls on both measures (BHC $-0.60$, AHC $-0.91$ points) and the Gini declines slightly ($-0.0028$). At 1 per cent, which raises \pounds 17.0 billion of net revenue, BHC poverty rises 0.83 points while AHC poverty is still 0.30 points below baseline. At higher rates poverty and measured income inequality rise substantially --- at 5 per cent, BHC poverty rises by 12.0 points and the Gini by 0.10 --- because the static analysis levies cash charges on asset-rich, income-poor households with no offsetting use of the additional revenue. These upper-grid rows are the gross incidence of the tax alone and not complete reform packages: any practical proposal at those rates would recycle the net revenue, and the rows should not be read as forecasts of what such a package would do. The wealth Gini (0.7038) is mechanically unchanged, since the simulation does not model asset-price responses; Section~\ref{sec:robustness} reports what capitalisation would imply for it.

\begin{table}[htbp]
\centering
\caption{Poverty and inequality by LVT rate (council tax abolished throughout)}
\label{tab:poverty}
\begin{tabular}{lrrrrr}
\toprule
Rate & Poverty BHC & $\Delta$ BHC & Poverty AHC & $\Delta$ AHC & $\Delta$ Gini \\
     & (\%) & (pp) & (\%) & (pp) & \\
\midrule
Baseline & 11.53 & ---      & 14.75 & ---      & --- \\
0.5\%    & 10.94 & $-$0.60  & 13.84 & $-$0.91  & $-$0.0028 \\
0.77\%   & 11.58 & $+$0.04  & 14.04 & $-$0.70  & $+$0.0003 \\
1.0\%    & 12.36 & $+$0.83  & 14.45 & $-$0.30  & $+$0.0033 \\
1.5\%    & 13.82 & $+$2.29  & 15.62 & $+$0.87  & $+$0.0108 \\
2.0\%    & 15.55 & $+$4.02  & 17.56 & $+$2.81  & $+$0.0197 \\
3.0\%    & 18.86 & $+$7.32  & 21.52 & $+$6.78  & $+$0.0417 \\
5.0\%    & 23.70 & $+$12.17 & 26.61 & $+$11.86 & $+$0.1013 \\
\bottomrule
\end{tabular}

\medskip
{\footnotesize\emph{Note:} Absolute poverty against baseline-fixed thresholds, before (BHC) and after (AHC) housing costs. Gini refers to equivalised household net income; baseline value 0.3641. Rates above 0.77\% raise net revenue that is not recycled in the simulation. \emph{Source:} Authors' calculations using PolicyEngine UK.}
\end{table}

\begin{figure}[htbp]
\centering
\safefig{../results/figures/fig8_poverty_gini_by_rate.png}
\caption{Poverty and Gini changes across the LVT rate grid.}
\label{fig:povgini}
\end{figure}

\subsection{Validation and robustness}
\label{sec:robustness}

\paragraph{The swap has no benefit interactions.} Running the reform through the full tax and benefit model and comparing with the arithmetic $\Delta y_i = \mathrm{CT}_i - r L_i$, we find the two agree to within \pounds 0.25 for every household at every rate tested (0.5, 1.0 and 2.0 per cent), and the share of households for which they differ by more than \pounds 1 is zero. Abolishing council tax and levying a land tax changes no means-tested entitlement in the UK system: council tax reduction is itself abolished with the tax, and neither instrument enters the income or capital tests for Universal Credit, Pension Credit or the legacy benefits. The consequence is that the microdata, not the tax-benefit calculator, do all the work in this reform --- the calculator's contribution is the council tax baseline, the equivalisation and the statutory poverty thresholds. We regard this as worth stating explicitly because it is easy to assume otherwise, and it means the results can be reproduced and audited without re-running the model.

\paragraph{Robustness.} Table~\ref{tab:robustness} varies the replacement target, the composition of the base and the instrument. The qualitative results are stable. The share of households gaining stays within a 65--70 per cent band across every specification. The budget-neutral rate ranges from 0.72 to 1.07 per cent, the top of that range being the household-land-only base, which is 28 per cent smaller. Calibrating the replacement target to OBR receipts rather than modelled liabilities lowers the rate to 0.72 per cent and, mechanically, makes every group better off --- it is a \pounds 3.9 billion net tax cut rather than a neutral swap, and we therefore treat it as a bound rather than a competing central case. Excluding Northern Ireland, whose households would otherwise pay the land tax on top of domestic rates, raises the rate to 0.79 per cent and leaves the distribution essentially unchanged. Assuming half of corporate equity is foreign-owned and therefore outside the household base raises the required rate to 0.90 per cent, the largest single sensitivity in the table, but still leaves 68.0 per cent of households gaining. Perturbing every regional land share by $\pm 10$ per cent, or rescaling household land to the un-uprated ONS 2024 benchmark, changes the rate by at most 0.06 percentage points and the decile effects by tens of pounds.

The final row is a comparator rather than a robustness check: a proportional tax on total property value, structures included, at the 0.73 per cent rate that raises the same revenue. It is less progressive on the wealth ranking --- the top wealth decile loses \pounds 4{,}217 rather than \pounds 5{,}752 --- and it produces fewer winners (64.4 against 68.2 per cent), because the base is spread more evenly across the property-owning population instead of being concentrated where location value is highest. That difference is the quantitative content of the land/structures distinction.

\begin{table}[htbp]
\centering
\footnotesize
\caption{Robustness: revenue-neutral rate and distributional summary under alternative assumptions}
\label{tab:robustness}
\begin{tabular}{lrrrrrrr}
\toprule
 & Base & Rate & Winners & Decile 1 & Decile 10 & Wealth 10 & $\Delta$ Poverty \\
Specification & (\pounds tn) & (\%) & (\%) & (\pounds) & (\pounds) & (\pounds) & AHC (pp) \\
\midrule
Central                              & 7.46 & 0.77 & 68.2 & $+$647 & $-$963    & $-$5{,}752 & $-$0.70 \\
OBR receipts target                  & 7.46 & 0.72 & 69.1 & $+$698 & $-$730    & $-$5{,}177 & $-$0.73 \\
Great Britain only                   & 7.34 & 0.79 & 67.9 & $+$639 & $-$1{,}016 & $-$5{,}747 & $-$0.56 \\
Household land only                  & 5.40 & 1.07 & 69.6 & $+$574 & $-$688    & $-$5{,}364 & $-$0.56 \\
Half of corporate land foreign-owned & 6.43 & 0.90 & 68.0 & $+$617 & $-$848    & $-$5{,}589 & $-$0.59 \\
Household land scaled to ONS 2024    & 7.10 & 0.81 & 68.2 & $+$651 & $-$977    & $-$5{,}772 & $-$0.71 \\
Land shares $-10$\%                  & 6.92 & 0.83 & 68.2 & $+$653 & $-$985    & $-$5{,}782 & $-$0.70 \\
Land shares $+10$\%                  & 8.00 & 0.72 & 68.2 & $+$642 & $-$945    & $-$5{,}726 & $-$0.56 \\
\midrule
\emph{Comparator:} proportional property tax & 7.91 & 0.73 & 64.4 & $+$605 & $-$674 & $-$4{,}217 & $-$0.76 \\
\bottomrule
\end{tabular}

\medskip
{\footnotesize\emph{Note:} Each row sets the rate so that revenue equals the replacement target on that row's base. ``Decile 1'' and ``Decile 10'' are average net income changes by income decile; ``Wealth 10'' is the top wealth decile. \emph{Source:} Authors' calculations using PolicyEngine UK; \texttt{analysis/extensions.py}.}
\end{table}

\paragraph{Capitalisation.} The central results are flow accounting and assume no capitalisation. If instead the levy is fully anticipated and permanent, it capitalises into land prices: a rate $r$ discounted at $\rho$ implies a proportional fall of $r/\rho$ in land values. At $\rho = 5$ per cent the 0.77 per cent rate implies a 15.4 per cent fall in land prices and an aggregate wealth loss of \pounds 1.15 trillion; at $\rho = 3$ per cent, a 25.7 per cent fall and \pounds 1.92 trillion. This is a transfer to the state of an asset value that was itself partly the capitalisation of past public investment, and under the textbook assumption it lands overwhelmingly on existing owners at announcement rather than on subsequent purchasers. Whether it does so in practice is contested: \citet{hoj2018} find full capitalisation in Danish data, whereas \citet{nielsson2024} cannot reject zero and argue the burden is shared with tenants and future buyers. We report the arithmetic under full capitalisation because it bounds the wealth effect, not because the evidence compels it. Table~\ref{tab:capitalisation} in the appendix reports the implied one-off wealth loss by wealth decile: at $\rho=3$ per cent the top wealth decile loses \pounds 280{,}331 of land value, 6.5 per cent of its total wealth, against \pounds 323 for decile 3. Two points follow. First, the wealth-decile gradient of the reform is steeper still on a stock than on a flow basis. Second, the flow results in Sections~\ref{sec:income_impact}--\ref{sec:wealth_impact} should be read as impact incidence in the first years, before the announcement effect on prices works through.
